\section{Generation along the dlt boundary}\label{sec:boundary}

We prove that the restrictions of a nef adjoint to the separate log
canonical strata fit together to generate its restriction to the whole
reduced boundary.  Semiampleness on the normal components does not by itself
supply compatible sections on their union.  The proof has three
parts.  A Mori contraction reduces the obstruction to restriction from a
stratum to one comparison of residues.  Crepant transport respects all
lower residues.  Finally a finiteness argument permits us to impose every
comparison simultaneously by taking products of sections.

Throughout this section, \(n\geq1\) and \(\Ggood_j\) is assumed for every
\(j<n\).  We use ordinary dlt pairs with effective rational boundary.
A \emph{stratum} is an irreducible log canonical center; we also allow the
ambient space itself when describing adjunction.  A \emph{boundary stratum}
is a stratum contained in the reduced floor.  An \emph{incidence} is an
inclusion of strata, and it is \emph{immediate} when the smaller stratum has
codimension one in the larger.  On a crepant SNC model, a \emph{unit
component} is a boundary component of coefficient one, and a \emph{unit
stratum} is an irreducible component of an intersection of unit components.
All residue comparisons are in a common divisible even adjoint degree.

\begin{theorem}[Generation on the reduced boundary]\label{floor:generation}
Assume \(\Ggood_j\) for every \(0\leq j<n\).
Let \((V,B)\) be a normal irreducible compact K\"ahler dlt \(n\)-fold with
effective rational boundary \(B\).  Suppose that \(V\) is globally
\(\Q\)-factorial, that \((V,B)\) satisfies the lc-strata resolution
convention of Definition~\ref{def:lc-strata-resolution}, and that the actual
\(\Q\)-Cartier adjoint \(J=K_V+B\) is analytically nef.
Then the restriction of an actual positive Cartier multiple of \(J\) to
the whole reduced space \(S=\lfloor B\rfloor\) is globally generated.
\end{theorem}

The assertion is vacuous if \(S=\varnothing\).  The remainder of the
section proves it when the floor is nonempty.  The construction of
compatible sections follows the admissible-section method of
Fujino~\cite[Section~4]{Fujino2000}; the main work here is to establish its
restriction and finiteness inputs for strata of arbitrary dimension in the
compact K\"ahler setting.

\subsection{Adjunction, the semiample systems, and their comparisons}

We first fix precisely which sections have to agree.  The following local
facts retain the actual line, including the residue identifications at
every intersection.

\begin{lemma}[Strata and descent of residues]\label{floor:adjunction}
For each stratum \(Z\subseteq V\) there is an effective rational boundary
\(B_Z\) on the normal compact K\"ahler space \(Z\) such that \((Z,B_Z)\)
is dlt and, in every sufficiently divisible even degree \(q\),
\begin{equation}\label{floor:residue-line}
 \mathcal O_Z\bigl(q(K_Z+B_Z)\bigr)
       \simeq \mathcal O_V(qJ)|_Z.
\end{equation}
This is the actual meromorphic iterated-residue identification.  The prime
components of \(C_Z:=\lfloor B_Z\rfloor\) are precisely the strata
\(D\subset Z\) in an immediate incidence.  Sections of \(\mathcal O_V(qJ)|_D\) on
these components \(D\) which have equal residues on every common lower
stratum descend uniquely to a section of the restricted line on the entire
reduced \(C_Z\).  The same statement applies to the components of
\(S\subseteq V\).
\end{lemma}

\begin{proof}
Choose one resolution \(p:\widehat V\to V\) witnessing
Definition~\ref{def:lc-strata-resolution}, and write
\(K_{\widehat V}+\widehat B=p^*(K_V+B)\) using the actual meromorphic
identification.  Its unit components are the strict transforms
\(\widehat S_i\) of the finitely many components \(S_i\) of \(S\).
Every component \(\widehat Z\) of an intersection of distinct
\(\widehat S_i\)'s maps birationally onto a stratum \(Z\), because the
map is an isomorphism at its general point.  Conversely every stratum is
obtained this way.  The general SNC locus makes both the component and its
index set unique.  Compactness makes the collection finite.
The empty intersection is \(\widehat V\) itself.  Write
\(p_Z=p|_{\widehat Z}:\widehat Z\to Z\), so \(p_V=p\).

Here are the adjunction and normality details, including the facts used at
deeper intersections.  For a chosen set of indices, keep their coefficients
one and allow each unused coefficient to be either one or \(1/2\).  Write
\(B^{\mathbf e}=B-\sum_i(1-e_i)S_i\).  Global \(\Q\)-factoriality makes
this an actual rational-line operation, and
\[
 \widehat B^{\mathbf e}=\widehat B-
                  p^*\sum_i(1-e_i)S_i.
\]
The subtracted divisor is effective.  All exceptional coefficients remain
below one, and the remaining unit components are exactly the retained
strict transforms.  Successive SNC adjunction on \(\widehat Z\) gives
\[
 (K_{\widehat V}+\widehat B^{\mathbf e})|_{\widehat Z}
       =K_{\widehat Z}+\widehat B_{\widehat Z}^{\mathbf e}.
\]
Here and below such a restriction equality means equality of the
meromorphic residue maps in degree \(q\), and hence equality of the
actual lines.  If two orders of residue differ by an ordering sign, its
\(q\)-th power is one.

We induct on the number of chosen indices.  At each stratum already
constructed, we retain normality, effectivity of every weighted different,
and the crepant SNC residue comparison on \(\widehat Z\).  For this chosen
comparison the exceptional coefficients are below one, the unit components are the retained strict
intersections, and the full-weight pair is dlt.  These assertions hold for
\(Z=V\).  Lowering unused unit components preserves dlt; lowering all of
them makes the effective pair at the current stratum klt.  It has rational
singularities and is Cohen--Macaulay by the dimension-free
floor-connectedness lemma \cite[Lemma~2.1]{BL}.  If just one unused
coefficient is retained, the unit divisor of the induced SNC boundary is a
disjoint union \(R\) of smooth divisors on \(\widehat Z\).  The same lemma gives
\[
 (p_Z)_*\mathcal O_R=\mathcal O_{p_Z(R),\mathrm{red}}.
\]
The fibers are connected, so the images of different connected components
of \(R\) cannot meet.  Each component maps birationally to its image; the
displayed equality, factored through finite normalization, makes that
image normal.  It is compact K\"ahler by restriction of the local K\"ahler
potentials from \(V\).

For every allowed weight choice, push forward the SNC different to define
the rational different on this new stratum.  In codimension one, the residue of a local frame of the
ambient invertible line identifies its divisorial adjoint with the
restricted line.  Reflexive extension on the normal stratum gives
Equation~\eqref{floor:residue-line} everywhere.  Pulling it back agrees with the
original SNC residue on a dense open and therefore everywhere, proving
crepancy as an actual meromorphic identification.

The different is effective.  At a general point of any of its primes,
the normal surface-slice calculation of \cite[Lemma~6.2]{BL}, which is
stated in every dimension, preserves exactly the residue order and reduces
it to adjunction on a normal surface germ with effective boundary and a
smooth marked curve.  For completeness, on a minimal resolution of that
surface, write the crepant boundary as the effective strict transform plus
an exceptional divisor \(E\).  For each exceptional curve \(D\),
\[
 E\cdot D=-(K_{\widetilde T}+B^{\rm str})\cdot D\leq0.
\]
Indeed \(B^{\rm str}\cdot D\geq0\), and adjunction gives
\(K_{\widetilde T}\cdot D=2p_a(D)-2-D^2\geq0\); minimality excludes a
smooth rational exceptional \((-1)\)-curve.  The negative-definite
exceptional intersection matrix gives \(E\geq0\): if \(E=P-N\) with
disjoint nonnegative parts and \(N\ne0\), then
\(E\cdot N=P\cdot N-N^2>0\), contrary to the preceding inequalities.
The strict transform is finite birational over the smooth marked curve,
and hence isomorphic to it.  Adjunction to that smooth strict transform
thus has nonnegative coefficient.  This is the original coefficient by the
surface-slice residue calculation.

The full-weight pair is dlt and has the asserted floor.  On an SNC model,
at the general point of the center of a divisorial valuation \(v\), take
normal parameters \(x_1,\ldots,x_a\), giving an unmarked parameter boundary
weight zero.  The SNC discrepancy inequality is
\[
 a(v;\widehat Z,\widehat B_{\widehat Z})
       \geq \sum_{i=1}^a(1-b_i)v(x_i).
\]
Every \(v(x_i)>0\).  Thus a zero-discrepancy center is an intersection
of unit components.  Those intersections meet the isomorphism locus,
while every exceptional component has coefficient below one.  To produce
a dlt resolution for the restricted map, principalize its nonisomorphism
locus and resolve together with the ordered SNC boundary.  That locus
contains no entire log canonical center, so the displayed discrepancy
inequality makes every new exceptional coefficient strictly below one.
This is the dlt resolution criterion.  It also identifies the unit primes
downstairs with the next incident strata.

Finally let \(R_{\rm all}\) be the whole unit union on \(\widehat Z\).  Sections
on its smooth components which agree on all their SNC intersections glue
by the elementary equalizer for the coordinate ideals of an SNC union.
Even iterated residues make the equalities independent of the chosen chain.
Apply the floor-connectedness lemma to the full-weight pair and this full
unit union.  It gives
\((p_Z)_*\mathcal O_{R_{\rm all}}=\mathcal O_{C_Z}\).
Projection formula therefore descends the glued section uniquely to the
restricted invertible line on the reduced \(C_Z\).  The proof with \(Z=V\)
gives the assertion for \(S\).
\end{proof}

Each boundary stratum has dimension less than \(n\).  The class of
\(J_Z:=K_Z+B_Z\) is nef, since it is the restriction of the nef adjoint
class.  Proposition~\ref{bir:resolution-transfer}, applied using
\(\Ggood_{\dim Z}\) on a log resolution, makes the actual rational line
\(J_Z\) semiample.  Increase \(q\) once
for all so that
\(L_Z:=\mathcal O_Z(qJ_Z)=\mathcal O_V(qJ)|_Z\) is generated for every
boundary stratum.  Its map and Stein factorization give
\begin{equation}\label{floor:semiample-map}
 f_Z:Z\longrightarrow Y_Z,\qquad L_Z=f_Z^*N_Z,
\end{equation}
where \(Y_Z\) is normal projective, \(f_Z\) has connected fibers, and
\(N_Z\) is ample.  These are actual line identities.  We call \(C_Z\)
\emph{vertical} if \(f_Z(C_Z)\ne Y_Z\); the empty floor is vertical.
If \(C_Z\) dominates \(Y_Z\), then for every \(k\geq1\) the restriction
\[
 H^0(Z,L_Z^k)\longrightarrow H^0(C_Z,L_Z^k|_{C_Z})
\]
is injective.  Indeed, every section comes from \(Y_Z\); if its pullback
vanishes on \(C_Z\), surjectivity makes it vanish at every point of \(Y_Z\).

Fix a K\"ahler class \(c\) on \(V\), and let \(c_Z\) be its restriction
to each stratum.  For boundary strata \(Z,Z'\) of the same dimension, an
\emph{arrow} is a proper bimeromorphic comparison with a projective
resolution whose source is smooth and compact K\"ahler,
\[
 \begin{tikzcd}[column sep=large]
 & T\arrow[dl,"p"']\arrow[dr,"p'"]&\\
 Z&&Z'
 \end{tikzcd}
\]
such that the two pulled-back adjoint lines are equal as invertible
subsheaves of meromorphic \(q\)-pluricanonical forms on \(T\), and
\begin{equation}\label{floor:class-comparison}
 p^*c_Z-(p')^*c_{Z'}\in\NS(T)_{\R}.
\end{equation}
Here \(\NS(T)_{\R}\) is the real span of first Chern classes of
holomorphic lines, with torsion removed.  The degree-two formula for a
smooth modification says that its new classes are exceptional divisor
classes.  It follows that Equation~\eqref{floor:class-comparison} is independent
of further resolution and is preserved by composition.  These arrows
therefore form a groupoid on the finite set of strata of each dimension.
An arrow \(\gamma:Z\dashrightarrow Z'\) transports sections by
\[
 (p')^*(\gamma_*s)=p^*s.
\]
Normality gives existence and uniqueness.  Transport commutes with
multiplication and identifies the complete systems and their Stein targets
in Equation~\eqref{floor:semiample-map}.
The equality of meromorphic lines defines this transport of adjoint sections; the class
congruence will control self-arrows when the floor is vertical.

\subsection{One residue comparison suffices for restriction}

The only obstruction to extending a section from \(C_Z\) is its possible
variation among the connected components of a general floor fiber.  We show
that all these components meet one general Mori fiber, whose floor is
connected or consists of two points.  Thus at most one residue comparison
is needed.

\begin{lemma}[Restriction and a Mori link]\label{floor:restriction}
For each positive-dimensional boundary stratum \(Z\), there is either no
comparison or one arrow between two components of \(C_Z\), allowing a
component to be compared with itself, with the following property.  In all
sufficiently large divisible degrees \(k\), a section of \(L_Z^k|_{C_Z}\)
extends to \(Z\) if its two residues agree under that arrow.  With no
arrow, every such section extends.  The degree bound is uniform over
sections.  When present, the arrow compares the two coefficient-one points
on general \(\PP^1\) fibers of a projective Mori contraction on a
bimeromorphic compact K\"ahler model of \(Z\).
\end{lemma}

\begin{proof}
We use the dimension-free torsion-free restriction lemma
\cite[Lemma~7.1]{BL}.  In its notation, for an effective compact K\"ahler
lc pair \((T,G)\), klt away from its reduced floor \(C\), and a connected
map \(f:T\to P\) to a normal compact base from which an actual positive
adjoint multiple is pulled back, it asserts that
\(R^1f_*\mathcal O_T(-C)\) is torsion-free.  The ideal sequence gives
\begin{equation}\label{floor:restriction-obstruction}
 \mathscr Q:=\coker(\mathcal O_P\longrightarrow f_*\mathcal O_C)
       \lhook\joinrel\longrightarrow R^1f_*\mathcal O_T(-C).
\end{equation}
For \(T=Z\) and \(P=Y_Z\), if \(C_Z\) is vertical, \(\mathscr Q\) has
proper support and is zero.  Serre vanishing for the kernel of
\(\mathcal O_P\to f_*\mathcal O_C\), tensored with \(N_Z^k\), extends
every restriction in a uniform large-degree tail.  If \(C_Z\) dominates,
the same display says that a section extends as soon as its values are
constant on the reduced floor fiber over a dense open of \(Y_Z\): its
class in the torsion-free \(\mathscr Q\otimes N_Z^k\) then vanishes, and
the local lifts glue.  They are unique because \(C_Z\) dominates \(Y_Z\),
so \(\mathcal O_{Y_Z}\to(f_Z)_*\mathcal O_{C_Z}\) is injective.
This uses neither reduced special scheme
fibers nor base change at their points.

Suppose now that \(C_Z\) dominates \(Y_Z\).  There is an effective rational
Cartier divisor \(D\) with \(\Supp D=\Supp C_Z\): restrict to \(Z\) the
globally \(\Q\)-Cartier floor primes of \(V\) not containing \(Z\).
For small rational
\(\epsilon>0\), \((Z,B_Z-\epsilon D)\) is effective klt.  Apply
Corollary~\ref{prog:small-model} to obtain a projective small strong
\(\Q\)-factorialization \(\rho:Z_0\to Z\), crepant also for the full pair.
Write \(\Delta_0\) for its full boundary and \(D_0=\rho^*D\).
The full adjoint
\[
 P=K_{Z_0}+\Delta_0
\]
is the actual semiample rational line pulled back from \(Y_Z\), and the
lowered klt adjoint is \(J_{\rm low}=P-\epsilon D_0\).  Put
\(d=\dim Z\), choose \(m>0\) with \(mP\) Cartier and generated, and take
\(b>4md\).  Consider
\[
 H=J_{\rm low}+bP=(1+b)P-\epsilon D_0.
\]
On a general fiber of the map to \(Y_Z\), its class is
\(-\epsilon\{D_0\}\).  This is the negative of a nonzero effective
divisor because the floor dominates \(Y_Z\), so its pairing with a
K\"ahler power is negative.  If \(H\) were pseudo-effective, a positive
current representing it would restrict to a positive current on almost
every resolved smooth fiber, contradicting this pairing.  Thus \(H\)
is not pseudo-effective.

Apply Lemma~\ref{prog:nef-trivial} to the lowered klt pair and the nef
rational line \(P\).  Its \(H\)-program with K\"ahler scaling has ordinary
\(J_{\rm low}\)-negative steps, is finite, and ends in a projective Mori
contraction.  Every step, including the final Mori ray, is \(P\)-trivial.
The same actual Cartier line \(mP\) descends at each birational step.
Consequently its semiample map to the fixed base \(Y_Z\) persists on
every working model.  That map is constant on each connected final Mori
fiber and hence factors through the Mori contraction.  We obtain
\[
 T\xrightarrow{u}W\xrightarrow{g}Y_Z.
\]
Both maps have connected fibers.  The full adjoint on each model is the
trace of \(P\).  On a common resolution of a birational step, its natural
meromorphic canonical comparison is an exceptional divisor with zero
class, since the two actual lines descend from the same line on the
contraction base.  Exceptional negativity applied to both signs makes
that divisor zero.  Thus the full pair remains crepant, including its
meromorphic adjoint identifications.  The full transformed pair
\((T,G)\) is effective lc and is klt away from its full floor \(C^+\):
there it agrees with the transformed lowered klt pair.  The program is
nonextracting, so the transformed divisor \(D^+\) has
\(\Supp D^+=\Supp C^+\).  Since the full adjoint is pulled back from
\(Y_Z\), the lowered adjoint on a \(u\)-fiber is
\(-\epsilon D^+\).  Hence \(D^+\) is relatively ample over \(W\).

Compare the connected components of a general floor fiber on the original
small model and on \(T\).  On a common projective log resolution the
full crepant subboundary is the same.  Its unit union maps to both floors
with connected fibers by \cite[Lemma~2.1]{BL}.  Proper closed surjections
with connected fibers preserve connected components, also after restricting
over a point of \(Y_Z\).  Thus the connected components of the two
underlying floor fibers correspond.

Every connected component of the floor fiber on \(T\) meets one common
general Mori fiber.  Indeed \(C^+\to W\) is surjective because \(C^+\)
is the support of the relatively ample \(D^+\).  Apply
Equation~\eqref{floor:restriction-obstruction} to \(u\).  Surjectivity of
the floor makes \(\mathcal O_W\to u_*\mathcal O_{C^+}\) injective; the
displayed torsion-free cokernel then makes \(u_*\mathcal O_{C^+}\)
torsion-free.  Let
\[
 C^+\longrightarrow\overline C\longrightarrow W
\]
be the Stein factorization.  The finite space \(\overline C\) is reduced,
and every irreducible component of
\(\overline C\) dominates \(W\).  To see the
last implication algebraically, a vertical minimal prime of a finite
reduced algebra over a domain would give, by prime avoidance, a nonzero
element killed by a nonzero element of the domain.  After shrinking \(Y_Z\)
to a dense open, the fiber-dimension theorem gives every component of
\(\overline C_y\) dimension \(\dim W-\dim Y_Z\).  The fiber \(W_y\) is irreducible of that
dimension: a resolution of \(W\) still has connected fibers over \(Y_Z\),
and a general one is smooth and connected and surjects onto \(W_y\).
Finiteness now makes each component of \(\overline C_y\) dominate \(W_y\).
Consequently each connected component of \(C^+_y\) meets \(u^{-1}(w)\)
for the same general \(w\in W_y\).

The reduced general fiber \(F\) of \(u\) is irreducible by the same
resolution argument.  If \(\dim F\geq2\), the support of the effective
ample Cartier multiple of \(D^+|_F\), whose support is \(C^+\cap F\),
is connected.  One elementary
proof, valid even if \(F_{\rm red}\) is not normal, cuts by general
hyperplanes to an integral projective surface.  If the ample divisor
split into disjoint nonzero Cartier parts \(D_1,D_2\), their pullbacks
to a resolution would be orthogonal, while the projection formula gives
\(v^*(D_1+D_2)\cdot v^*D_i>0\) for both \(i\).  Thus both
\((v^*D_i)^2>0\), contradicting the surface Hodge index theorem.
If \(\dim F=1\), choose the fiber also off the singular, boundary
intersection, and ramification loci.  It is a smooth connected curve and
\[
 0=\deg(K_F+G|_F)=2g(F)-2+\deg(G|_F).
\]
There is a unit point, so \(F\simeq\PP^1\).  Its floor has one or two
points; if there are two they exhaust the horizontal boundary.  In every
connected case the preceding common-fiber observation makes the general
floor fiber connected, so no comparison is needed.  In the remaining case
the two unit points represent all its connected components.

Normalize the horizontal floor primes and take their Stein factorizations
over \(W\).  If there are two primes, each finite Stein map has degree one
and is an isomorphism over normal \(W\); their main fiber product gives
a proper bimeromorphic comparison.  If there is one prime, its normal Stein
space is a double cover of \(W\).  The exchange on its general fiber
extends over the branch locus: the reduced horizontal non-diagonal
component of the double fiber product has finite birational projections
to the normal cover and hence both projections are isomorphisms.  Its graph
lifts to the normalized prime.  Composing with the program graphs transports
the comparison to two original normal floor primes, possibly the same one.
All graph projections are projective, and the primes survive birationally
because the program extracts no divisors.

This comparison is crepant as a meromorphic residue identity.  On a smooth
ruled open, order the markings after an \'etale local cover, choose a base
volume \(\xi\), and put the markings at \(z=0,\infty\).  A pulled-back
\(q\)-pluriadjoint frame has the form
\[
 a(w)(dz/z)^{\otimes q}\otimes\xi^{\otimes q}.
\]
Its two residues are \(a(w)\xi^{\otimes q}\) and
\((-1)^q a(w)\xi^{\otimes q}\), which agree.  This is the two-marking
Poincar\'e-residue comparison of
\cite[Section~3, Definition~13 and Proposition~14]{Kollar2013Sources}.
On a resolved graph both restricted actual lines are pulled back from
\(Y_Z\); their meromorphic embeddings agree on this dense open and hence
everywhere, including along exceptional primes.  Equal residues therefore
give a constant value on the entire general floor fiber.  The horizontal
consequence of Equation~\eqref{floor:restriction-obstruction} extends the section.
Normality descends it from the small model.  Its restriction is the
prescribed section on the original reduced floor, since every original
floor prime is covered birationally and the two restrictions agree after
that pullback.

It remains to check Equation~\eqref{floor:class-comparison}.  Let \(\widehat T\)
be a common smooth resolution of the parent and Mori models and let
\(a,b:H\to\widehat T\) be the two maps from a resolved branch graph.
The resolution of the Mori model is an isomorphism near a general entire
\(\PP^1\) fiber: its centers have codimension at least two, hence cannot
dominate the ruling base.  Every holomorphic two-form on \(\widehat T\)
comes from the base on that ruled open.  In fact its vertical one-form
terms vanish on \(\PP^1\), and its remaining coefficients are constant
on that compact fiber.  Thus \(a^*-b^*\) kills \(H^{2,0}\), and also
\(H^{0,2}\) by conjugation.  This rational Hodge map on degree two has
image of type \((1,1)\).  The Lefschetz \((1,1)\) theorem puts its real
image in \(\NS(H)_{\R}\).  Apply this to the pulled-back parent class
\(c_Z\).  The link is therefore an arrow as defined above.
\end{proof}

\subsection{Transport respects every lower residue}

An arrow can contract a divisor to a higher-codimension stratum.  The next
local observation explains how to compare residues even in that case.
All valuations in its proof are local monomial valuations in an SNC chart.

\begin{lemma}[Unit strata above an SNC stratum]\label{floor:unit-strata}
Let \(p:T\to Z\) be a projective log resolution of a dlt pair, with
crepant SNC subboundary on the smooth \(T\).  Let \(R\subseteq T\) be a
unit stratum and let \(D=p(R)\) be its center.  Then
\(D\) is a dlt stratum.  There is a unit stratum \(R'\subseteq R\) such
that \(p(R')=D\) and \(R'\to D\) is birational.  In the iterated residue
comparison along \(R'\), the normal logarithmic Jacobian is \(\pm1\).
\end{lemma}

\begin{proof}
The center is a dlt stratum by the SNC discrepancy calculation in the
proof of Lemma~\ref{floor:adjunction}.  Work over the generic SNC part of \(D\),
of codimension \(a\), with normal coordinates \(x_1,\ldots,x_a\) for
the unit divisors.  The zero-discrepancy divisorial valuations centered
there are exactly the primitive rational monomial rays in
\(\R_{>0}^a\).  Here is a local verification that also fixes their
normalization.  The discrepancy inequality in the proof of
Lemma~\ref{floor:adjunction} says on every log-smooth model that the center
of a zero-discrepancy valuation is exactly a unit stratum: an additional
unmarked normal parameter, or a parameter with coefficient below one, would
give positive discrepancy.  Blow up that stratum.  If the positive normal
orders of the valuation are \(w_1,\ldots,w_b\), the chart indexed by a
smallest order \(w_k\) replaces these by \(w_k\) and the positive members
of \(w_i-w_k\).  The exceptional divisor is again unit, and the center
is exactly the new unit stratum.  While there is more than one positive
order their sum decreases.  Eventually the center is the generic point of
one unit divisor, where a normalized divisorial valuation has order one.
Reversing these ordinary toric blowups proves monomiality and primitive
normalization.  Conversely this subtraction algorithm for any primitive
positive integer vector terminates at order one, and reversing it extracts
that monomial valuation.  Comparison on a common graph gives uniqueness
of the valuation with those weights and of its center on any model.

For a unit component \(E\) meeting the inverse image of the generic SNC
part of \(D\), put
\[
 v_E=(\ord_E x_1,\ldots,\ord_E x_a)\in\Z_{\geq0}^a.
\]
If the unit components through a unit stratum \(Q\) meeting that inverse image are
\(E_1,\ldots,E_b\), define
\[
 \phi_Q:\R_{\geq0}^b\longrightarrow\R_{\geq0}^a,\qquad
 (t_1,\ldots,t_b)\longmapsto\sum_{j=1}^b t_jv_{E_j},
 \qquad \sigma_Q=\phi_Q(\R_{\geq0}^b).
\]
The source carries its standard lattice \(\Z^b\).  For the fixed
resolution \(T\), these cones subdivide the orthant over its interior.
On rational rays, \(\phi_Q\) sends a
monomial valuation on \(T\) to the same normalized valuation on \(Z\),
so it is injective on each cone.  Interiors of different cones are
disjoint, because the center and the weights of a monomial valuation on
\(T\) are unique.  Every rational ray in the orthant occurs: view its
valuation on \(T\) and use the same SNC description there.  There are
only finitely many cones over the generic part of \(D\); density of
rational rays therefore gives coverage.  Taking subsets of unit components
gives the faces, and uniqueness of valuations also on faces identifies
intersections as common faces.  Consequently \(\sigma_R\) is a face
of an \(a\)-dimensional cone \(\sigma_{R'}\) with the same downstairs
center, for a unit stratum \(R'\subseteq R\).

Let \(E_1,\ldots,E_a\) be the unit components through \(R'\).
The integer matrix \(A=(\ord_{E_j}x_i)\) of this full cone is unimodular.
Indeed every positive
primitive integral vector in the source cone is the normalized vector of
a divisorial valuation, so its image is primitive.  If a prime divided
\(\det A\), a nonzero kernel vector modulo that prime, lifted to a
positive primitive integral vector, would have nonprimitive image.  Hence
\(\det A=\pm1\).  The strata \(R'\) and \(D\) have the same dimension.
We next prove that their generically finite map has degree one.  In local coordinates at a
general point of \(R'\),
\[
 x_i=u_i(y,z)\prod_{j=1}^a y_j^{A_{ij}},
\]
where the \(u_i\) are units, the \(y_j\) are normal coordinates, and
\(z\) are coordinates along the stratum.  The tangential map is
generically \'etale in characteristic zero.  Since \(A^{-1}\) is integral,
replace \(y_j\) by
\(\widetilde y_j=y_j\prod_i u_i^{(A^{-1})_{ji}}\).  Then
\(x_i=\prod_j\widetilde y_j^{A_{ij}}\), and the map consisting of
\(\widetilde y\) and the downstairs tangential coordinates is locally
biholomorphic at a general \'etale point of \(R'\to D\).  Two distinct
points over a general point of \(D\) would therefore give two lifts of
nearby target points whose normals approach with any fixed positive rate
vector in the interior of this cone.  Varying the nonzero leading
coefficients fills an open set of nearby torus points, so such a target can
be chosen in the locus where \(p\) is an isomorphism.  This contradiction
proves that \(R'\to D\) is birational.  Finally, after taking residue
along \(R'\), the normal coefficient of
\(\bigwedge_i dx_i/x_i\) is \(\det A\) times that of
\(\bigwedge_jdy_j/y_j\); terms differentiating the units vanish in that
residue.  Since \(\det A=\pm1\), the even pluriresidue removes the sign.
\end{proof}

\begin{lemma}[Preservation of lower restrictions]\label{floor:transport}
Fix \(0\leq d<n\) and \(k\geq1\).  For each boundary stratum \(Z\) of
dimension at most \(d\), let \(s_Z\in H^0(Z,L_Z^k)\).  Suppose these
sections agree under the residue restriction for every incidence, and that
every arrow in dimensions below \(d\) carries its source section to its
target section.  For any \(d\)-arrow \(\gamma:Z\dashrightarrow Z'\) and
every component \(D'\) of \(C_{Z'}\),
\[
 (\gamma_*s_Z)|_{D'}=s_{D'}.
\]
Moreover verticality of \(C_Z\) is invariant under \(d\)-arrows.
\end{lemma}

\begin{proof}
The assertions are immediate for \(d=0\), when every floor is empty.
For \(d>0\), choose a common log resolution
\(p:T\to Z\), \(p':T\to Z'\) of the arrow.  For a target floor prime
\(D'\), start with its strict transform \(R\).  It is a unit stratum and
maps birationally to \(D'\).  Its source center \(p(R)\) is a stratum by
Lemma~\ref{floor:unit-strata}.  If \(R\to p(R)\) is not birational, that
lemma replaces \(R\) by a unit stratum inside it which maps birationally
onto the same source center.  The target center may then shrink.  If the
map to that center is not birational, apply the lemma on the target side,
and continue alternately.  Each necessary replacement strictly decreases
dimension.  The process therefore ends with one unit stratum mapping
birationally to lower strata \(D_0,D_0'\) on the two sides.

The image in the common system base remains unchanged throughout.  Indeed
the two maps from the resolved graph to the Stein targets are identified
by the arrow, and a replacement keeps its entire image on the side then
being treated.  Initially the image was \(f_{Z'}(D')\).  The terminal
comparison of \(D_0\) and \(D_0'\) is an arrow: restricting the actual
meromorphic identification to iterated residues gives crepancy, with the
logarithmic determinant sign removed by Lemma~\ref{floor:unit-strata};
restricting the line Chern classes in Equation~\eqref{floor:class-comparison}
gives the same condition modulo \(\NS_{\R}\).  Lower invariance thus
identifies the transported and assigned residues on \(D_0'\).

This equality determines the entire section on \(D'\).  The Stein map
of \(L_{D'}=L_{Z'}|_{D'}\) is the Stein factor of
\(f_{Z'}|_{D'}\), because the line induced on that Stein space is ample;
its target is finite over \(f_{Z'}(D')\).  The image of \(D_0'\) in this
target is a closed irreducible subset of full dimension, since it still
maps onto \(f_{Z'}(D')\), and therefore is the whole target.  Both
sections on \(D'\) come from it.  Equality after restriction to
\(D_0'\) proves equality on \(D'\).

The same construction shows that the image of every target floor prime
in the system base is the image of a lower source stratum, and the reverse
statement follows from the inverse arrow.  A floor dominates the system
base on one side exactly when it does on the other.  This proves the last
assertion.
\end{proof}

\subsection{Finite actions on vertical strata}

Once lower-dimensional sections are compatible and invariant,
Lemma~\ref{floor:transport} preserves their restrictions under every arrow.
If \(C_Z\) dominates \(Y_Z\), injectivity of restriction forces any such
transported extension to equal the assigned extension.  Thus only vertical
strata require the following argument for finite image.

\begin{proposition}[Finite image of self-arrows]\label{floor:finite-action}
Let \(Z\) be a boundary stratum with \(C_Z\) vertical.  For all sufficiently
large divisible \(k\), the self-arrows of \(Z\) have finite image on
\(H^0(Z,L_Z^k)\).
\end{proposition}

For \(m=qk\), a pluricanonical eigenform \(s\in H^0(Z,L_Z^k)\) with
\(\int_Z|s|^{2/m}<\infty\) becomes the \(m\)-th power of a holomorphic top
form on a smooth resolution of a cyclic root cover, so the auxiliary proposition below makes its
eigenvalue a root of unity.  A uniform bound on the middle Betti numbers of
selected resolutions of these covers in the fixed degree will then bound
the possible orders; boundedness of the representation will give finite
image. The auxiliary proposition uses the polarization congruence in
Equation~\eqref{floor:class-comparison}.

\subsubsection{Top-form scalars}

\begin{proposition}[A polarized top-form scalar]\label{floor:top-form-scalar}
Let \(U\) be a smooth connected compact K\"ahler manifold of dimension
\(d<n\), and let \(g:U\dashrightarrow U\) be bimeromorphic.  Suppose
that on a common smooth resolution the two pullbacks of some K\"ahler
class on \(U\) differ by an element of \(\NS_\R\).  If
\(0\ne\alpha\in H^0(U,K_U)\) and \(g^*\alpha=\mu\alpha\), then
\(\mu\) is a root of unity.
\end{proposition}

We first control the canonical Iitaka base; this step uses \(\Ggood_d\)
and the top form, with no polarization congruence.  To specify that base,
let \(U\) be a smooth connected compact K\"ahler \(d\)-fold with \(d<n\)
and \(\kappa(U,K_U)\geq0\).  A pluricanonical section makes \(K_U\)
pseudo-effective, so \(\Ggood_d\) gives a smooth
modification \(\pi:U'\to U\) and
\[
 \pi^*K_U\sim_{\Q}P+E_U,\qquad
 P\ \text{semiample},\qquad E_U=N(\pi^*K_U).
\]
Choose \(\ell>0\) sufficiently divisible that the connected semiample map
and its Stein target satisfy
\[
 f:U'\longrightarrow Y,\qquad
 \mathcal O_{U'}(\ell P)=f^*N,
\]
where \(Y\) is normal projective and \(N\) is ample, and that
Lemma~\ref{bir:sections} identifies every graded piece in
\[
 R_\ell(U):=\bigoplus_{j\geq0}H^0(U,j\ell K_U)
 \simeq \bigoplus_{j\geq0}H^0(Y,N^j).
\]
This is the full section ring of the ample line \(N\), so
\(Y\simeq\Proj R_\ell(U)\).  Passing to a further divisible Veronese gives
the same Proj.  We call \(Y\) the \emph{canonical Iitaka base}.  A
bimeromorphic self-map of \(U\) acts on this graded ring and therefore
induces an automorphism of \(Y\).

\begin{lemma}[Finite action on the canonical Iitaka base]\label{floor:finite-base}
Let \(U\) be a smooth connected compact K\"ahler manifold of dimension
\(d<n\), let \(0\ne\alpha\in H^0(U,K_U)\), and let \(g:U\dashrightarrow U\)
be bimeromorphic with \(g^*\alpha=\mu\alpha\).  The induced automorphism
\(h\) of the canonical Iitaka base \(Y=\Proj R_\ell(U)\), for a sufficiently
divisible \(\ell\) as above, has finite order.
\end{lemma}

\begin{proof}
The top form gives \(\kappa(U,K_U)\geq0\), so the preceding construction
applies.  The assertion is
immediate for a point base.  Otherwise use
\(f:U'\to Y\) above, and write \(\alpha\) also for its pullback to \(U'\).
For each positive integer \(k\) the
canonical integral \(s\mapsto\int_U|s|^{2/k}\) on \(H^0(U,kK_U)\) is
positive, continuous, and invariant under bimeromorphic change of variables.
Choose \(j\) such that \(N^j\) is very ample.  The powers of \(g\) and
their inverses are bounded on the corresponding degree \(j\ell\), which
defines \(Y\).  This linear action is semisimple with eigenvalues of
absolute value one.  Also \(|\mu|=1\), by integration of
\(\alpha\overline\alpha\).  Hence the finite measure
\[
 \nu=f_*(i^{d^2}\alpha\wedge\overline\alpha)
\]
is invariant under \(h\); exceptional sets have measure zero here.

Suppose that \(h\) has infinite order.  In the projective linear group of
this degree, the algebraic closure of a power of \(h\) is
then a positive-dimensional torus \(T\), acting faithfully on \(Y\).
Indeed a bounded linear operator is diagonalizable, and the connected
algebraic closure of a cyclic diagonal group is a torus.  We will obtain a
contradiction from any nontrivial one-parameter subgroup of \(T\).

Let \(Y^\circ\) be a dense Zariski open in the smooth locus where \(f\)
is smooth and the relative top form obtained from \(\alpha\) is nonzero.
Put \(r=\dim U'-\dim Y\).  A general fiber \(F\) of an Iitaka map has
\(\kappa(F)=0\).  It has a nonzero holomorphic top form here, so
\(h^0(F,K_F)=1\): two such forms would have a nonconstant ratio and give
positive Iitaka dimension.  If \(\xi\) is a local holomorphic volume
frame on \(Y^\circ\) and \(\sigma=\alpha/\xi\) is the relative top
form, fiber integration writes the measure as
\begin{equation}\label{floor:hodge-density}
 \nu=\|\sigma\|_{\rm Hdg}^{2}\,i^{(\dim Y)^2}\xi\wedge\overline\xi.
\end{equation}
The coefficient is smooth and strictly positive after this shrink.  We
include \(r=0\), when the top Hodge line is the one-dimensional degree-zero
line.

Set
\[
 Y^*=\bigcup_{j\in\Z}h^j(Y^\circ).
\]
Invariance of \(\nu\) under \(h\) makes the smooth positive densities
agree on overlaps, extending the density in
Equation~\eqref{floor:hodge-density} to \(Y^*\).
This open is invariant under \(T\).  Its complement is the
intersection of the closed algebraic sets \(h^j(Y\setminus Y^\circ)\),
which is closed algebraic by Noetherianity.  Its algebraic stabilizer
contains the cyclic group generated by \(h\), and therefore contains its
algebraic closure.  Choose a one-parameter subgroup \(\C^*\subset T\)
which acts nontrivially, and a general whole orbit in \(Y^*\).
The \(T\)-invariance of the domain supplies this whole orbit.  We now
construct a horizontal period metric on \(Y^*\) and prove that its
restriction to the orbit vanishes.

For a smooth K\"ahler family a total K\"ahler class polarizes the primitive
real variation in middle cohomology.  On a simply connected open of
\(Y^\circ\), take the smallest real subvariation of the full middle
cohomology containing the line \(F^r\) of top forms and its conjugate.
It is the intersection of the flat real Hodge subspaces with this property;
finite dimensionality reduces the intersection to a finite one.  This
construction commutes with restriction and continuation: in a flat
trivialization the condition of being a Hodge subspace, and of containing
the indicated line, is a real analytic equality for the Hodge projections,
so equality on one open continues on the simply connected cover.  The
primitive variation is one such subvariation, so the smallest one lies in
primitive middle cohomology.  Its polarization is the middle cup pairing
with the usual sign, independent of the total K\"ahler class, because no
power of that class occurs in the middle primitive pairing.  The horizontal
period metric of this polarized real variation is therefore intrinsic.

This possibly degenerate metric extends consistently to \(Y^*\).
To compare the metrics on an overlap of two translates of \(Y^\circ\),
shrink the overlap further so that a resolution of the fiberwise birational
graph is a smooth family over it.  The two pullbacks on middle cohomology are
flat Hodge embeddings into the cohomology of this common resolution and
carry the top lines to the same line.  The smallest subvariation there is
the pullback of the smallest subvariation on either side: it is contained
in both images, and applying the inverse embedding gives the reverse
inclusion.  Projection formula identifies their middle cup polarizations.
Their horizontal metrics therefore agree on this smaller open and, by
continuity, on the overlap.

We claim that the period metric restricts to zero on every whole orbit
of this one-parameter subgroup in \(Y^*\); constant orbits are immediate.
For a nonconstant orbit, pull it back under
\(\exp:\C\to\C^*\).  The negative upper bound for holomorphic sectional
curvature in horizontal directions of a period domain, together with the
curvature inequality for a pulled-back curve metric, applies to its local
period representations; it holds for real polarizations as well
\cite[Theorem~9.1]{GriffithsSchmid1969}.  On a disk of radius \(R\)
centered at an arbitrary point, translate the center to \(z=0\) and
write the pulled-back metric as \(\chi(z)|dz|^2\), continuous and smooth
where positive, with curvature
\(-2\chi^{-1}\partial_z\partial_{\bar z}\log\chi\le-\kappa<0\).
Compare it with
\[
 \chi_R(z)=\frac{4R^2}{\kappa(R^2-|z|^2)^2}.
\]
Direct calculation gives
\(\partial_z\partial_{\bar z}\log\chi_R=\kappa\chi_R/2\).
If the maximum of \(\chi/\chi_R\) exceeded one, it would occur where
\(\chi>0\) in the interior, since \(\chi_R\) diverges at the boundary
and \(\chi\) is bounded on the closed disk.  At that point,
\[
 0\ge\partial_z\partial_{\bar z}\log(\chi/\chi_R)
   \ge\frac{\kappa}{2}(\chi-\chi_R)>0,
\]
a contradiction.  Thus the Ahlfors--Schwarz comparison gives
\(\chi(0)\le4/(\kappa R^2)\) \cite{Ahlfors1938}.
Letting \(R\to\infty\) forces the metric to vanish
on \(\C\).  The argument is local and allows zeros of the
induced metric, so the local representations just constructed suffice.
Consequently their period maps are locally constant along each such orbit.
In any holomorphic volume frame the logarithm of the coefficient in
Equation~\eqref{floor:hodge-density} is harmonic along each orbit: the top Hodge line is
flat there, and \(\sigma\) is a nonzero holomorphic multiple of a flat
generator.

This is incompatible with finite mass.  Diagonalize the one-parameter
action in projective coordinates, and choose two coordinates nonzero at a
general point with different weights, of difference \(a\ne0\).  A small
smooth transverse slice \(Q\) on which their ratio is one gives a
holomorphic local biholomorphism onto its image
\[
 \Phi:\C^*\times Q\longrightarrow Y^*
\]
with fibers of size at most \(|a|\): the ratio at \(\Phi(z,t)\) is
\(z^a\).  The entire \(\C^*\) factor is allowed because \(Y^*\) is
invariant.  In the product frame \((dz/z)\wedge dt_1\wedge\cdots\wedge
dt_{\dim Q}\), write \(\Phi^*\nu\) with coefficient \(\rho(z,t)>0\).
For fixed \(t\), \(\log\rho(z,t)\) is harmonic on \(\C^*\).  Its
circular mean at \(|z|=e^u\) is \(A(t)u+B(t)\).  Jensen's inequality
therefore gives, with an irrelevant positive normalization,
\[
 \int_{\C^*}\rho(z,t)\frac{i\,dz\wedge d\overline z}{|z|^2}
 \ \geq\ c\int_{\R}e^{A(t)u+B(t)}\,du=\infty.
\]
Tonelli's theorem contradicts the finite mass of \(\Phi^*\nu\), which
is at most \(|a|\nu(Y)\) by the bounded covering degree.  This proves
that \(h\) has finite order.
\end{proof}

After a power, the scalar problem can thus be restricted to a fiber with
Kodaira dimension zero.  The polarization congruence then passes to that
fiber; the next two lattice arguments control its nonprojective factors.
In the next lemma, the hypothesis
on the origin of the model is what provides the degree-two exceptional
splitting; rational singularities alone are not used for that splitting.

\begin{lemma}[The scalar on a symplectic factor]\label{floor:symplectic-scalar}
Let \(Y\) be a terminal compact K\"ahler primitive symplectic space of
dimension \(2e<n\), obtained from a smooth model by the program of
Corollary~\ref{prog:terminal-kappa-zero}.  Let \(g:Y\dashrightarrow Y\) be
bimeromorphic and let \(0\ne\eta\in H^0(Y,\Omega_Y^{[2]})\).  Suppose
there is a class \(c\in H^{1,1}(Y,\R)\) whose two pullbacks on a common
resolved graph of \(g\) differ by real line Chern classes, and such that
on some projective resolution \(p:R\to Y\) with smooth compact K\"ahler source
\[
 p^*c=[\beta]+\{E\},
\]
where \(E\) is an exceptional real divisor and \(\beta\) is a smooth
closed semipositive \((1,1)\)-form, strictly positive on a nonempty open.
Then the scalar in \(g^*\eta=\lambda\eta\) is a root of unity.
\end{lemma}

\begin{proof}
We use the pure Hodge structure on \(H^2(Y,\Q)\) and its rational
Beauville--Bogomolov form \(q_Y\), up to positive rational scale.  This
form is positive on the real symplectic plane and Lorentzian on
\(H^{1,1}(Y,\R)\); see \cite[Lemma~2.1, Definition~5.2, and
Lemmas~5.3 and~5.7]{BakkerLehn}.  The pullback and \((1,1)\)
identifications for these rational singularities, extension of the
two-form \cite[Corollary~1.8]{KebekusSchnell2021}, and
Lemma~\ref{prog:cohomology-transport} give
\begin{equation}\label{floor:symplectic-splitting}
 H^2(R,\Q)=p^*H^2(Y,\Q)\oplus
       \langle[p\text{-exceptional primes}]\rangle_{\Q}.
\end{equation}
Indeed that program lemma gives the real \((1,1)\) equality with a direct
exceptional sum; the \((2,0)\) and \((0,2)\) summands compare by form
extension, and all the summands in the resulting equality are rational.

Let \(\mathcal E_R\) be the rational span of the \(p\)-exceptional prime
classes.  Identify \(H^2(R,\Q)/\mathcal E_R\) with \(H^2(Y,\Q)\) by
Equation~\eqref{floor:symplectic-splitting}, and define the lattice
\[
 \Lambda_Y=
 \im\bigl(H^2(R,\Z)/\mathrm{tors}\longrightarrow
                 H^2(R,\Q)/\mathcal E_R\bigr)
 \subset H^2(Y,\Q).
\]
This amounts to quotienting the integral lattice by its saturated
exceptional sublattice.  The lattice is independent of a refinement
\(v:R'\to R\): integral Gysin satisfies \(v_*v^*=1\), and the smooth
degree-two modification formula says that
\(x'-v^*v_*x'\) is exceptional and that the new exceptional span is
\(v^*\) of the old one plus the \(v\)-exceptional span.  A common
smooth refinement therefore identifies the quotient lattices from any two
resolutions over \(Y\).

The form \(\eta^e\) trivializes \(K_Y\).  On a resolution \(R_g\) of
the graph of \(g\), with projections \(p_1,p_2\), the divisors of the
two pulled-back volume forms agree.  Terminality says that their positive
components are exactly the exceptional primes.  Thus the exceptional lists
for the two projections are identical; in particular \(g\) is small in
both directions.  Let \(\mathcal E\) be their common exceptional rational
span, and let
\[
 \overline p_i^*:H^2(Y,\Q)\xrightarrow{\sim}H^2(R_g,\Q)/\mathcal E.
\]
Both maps carry \(\Lambda_Y\) onto the image of integral
graph cohomology in this quotient.  Hence
\(G=(\overline p_1^*)^{-1}\overline p_2^*\) is a rational Hodge
automorphism of \(H^2(Y,\Q)\) preserving \(\Lambda_Y\).

It also preserves \(q_Y\).  The two symplectic forms pull back up to the scalar
\(\lambda\), and integration of their top powers gives \(|\lambda|=1\).
Use the usual positive normalization of
\(W=(p_1^*\eta\,p_1^*\overline\eta)^{e-1}
=(p_2^*\eta\,p_2^*\overline\eta)^{e-1}\).  For \((1,1)\) classes
\(x,y\), the differences between \(p_1^*Gx\) and \(p_2^*x\) are
exceptional for both projections.  Expand
\(p_1^*Gx\,p_1^*Gy-p_2^*x\,p_2^*y\) into two terms, each with one such
exceptional factor and one factor pulled back by the appropriate projection.
Projection formula makes both pairings with \(W\) zero.  This is precisely
the \((1,1)\) part of the Beauville--Bogomolov formula.  The symplectic
plane is preserved isometrically because \(|\lambda|=1\), and the Hodge
decomposition is orthogonal.  Thus \(G\) is a \(q_Y\)-isometry.

The class \(c\) in the statement has \(q_Y(c)>0\). Since \(\beta\)
is nef over \(Y\), exceptional negativity applied to the divisor \(E\),
with \(\{E\}=p^*c-[\beta]\), gives \(E=\sum_i e_iE_i\geq0\). Pair with the positive
normalization \(W_R=(p^*\eta\,p^*\overline\eta)^{e-1}\).  Exceptional
classes pair to zero with \(p^*c\) by projection formula, and hence for
some constant \(b_e>0\)
\begin{equation}\label{floor:positive-bb}
 q_Y(c)=b_e\left(\int_R\beta^2 W_R+\sum_i e_i\int_{E_i}\beta W_R\right)>0.
\end{equation}
Every term is nonnegative.  The first is strictly positive on the common
open where \(\beta\) is positive and the symplectic form is nondegenerate.

Set \(H_\Q=H^2(Y,\Q)\) and
\(D_\Q=H_\Q\cap H^{1,1}(Y,\C)\).  Lefschetz \((1,1)\) on \(R\), together
with the line-trace assertion of Lemma~\ref{prog:cohomology-transport},
identifies \(D_\Q\) with the rational span of line Chern classes.
It is \(G\)-invariant because \(G\) is a rational Hodge automorphism.
The same line-trace assertion applied to the graph congruence gives
\(Gc-c\in D_\R\).  The line \(\C\eta\) survives in \(H_\C/D_\C\);
our immediate aim is to prove that all eigenvalues on this quotient have
absolute value one.  Write \(\Sigma=(H^{2,0}\oplus H^{0,2})_\R\), the
positive real two-plane.  There are three possibilities for the restriction
of the Lorentz form to \(D_\R\).

If \(D_\R\) is negative definite, its orthogonal projection \(c_0\) of
\(c\) to \(D_\R^\perp\) is fixed by \(G\) and has positive square by
Equation~\eqref{floor:positive-bb}.  On \(D_\R^\perp\), the invariant space
\(\Sigma\oplus\R c_0\) is positive definite and its complement is
negative definite.  All eigenvalues there have absolute value one.  If
\(D_\R\) contains a positive vector, it is nondegenerate with one
positive direction.  Its orthogonal complement is the positive plane
\(\Sigma\) plus a negative-definite space, giving the same conclusion.
In either case \(H_\R/D_\R\simeq D_\R^\perp\).

In the remaining case \(D_\R\) is negative semidefinite with radical a
rational isotropic line \(\ell\).  The lattice action on \(\ell\) is
\(\pm1\).  On \(\ell^\perp/\ell\) the form is the positive plane
\(\Sigma\) plus a negative-definite \((1,1)\) space, and the quotient
\(H_\R/\ell^\perp\) is dual to \(\ell\).  The invariant flag
\(0\subset\ell\subset\ell^\perp\subset H_\R\) therefore shows that
all eigenvalues on \(H_\R\) have absolute value one; this argument also
allows unipotent parts.  In all cases the eigenvalues on
\(H_\C/D_\C\) have absolute value one.  The rational quotient
\(H_\Q/D_\Q\) carries the preserved lattice
\[
 \overline\Lambda_Y=\Lambda_Y/(\Lambda_Y\cap D_\Q),
\]
where the intersection is saturated.  Thus
\(\lambda\) is an algebraic integer and all its algebraic conjugates have
absolute value one.  Kronecker's theorem makes it a root of unity.
\end{proof}

\begin{lemma}[The scalar on a torus]\label{floor:torus-scalar}
Let \(A\) be the integral Hodge automorphism of \(H^1(T,\Z)\) induced
by an automorphism of a compact complex torus \(T\).  Suppose a positive
Hermitian class \(c\in H^{1,1}(T,\R)\) satisfies
\((\bigwedge^2 A)c-c\in\NS(T)_\R\).  Then
\(\det(A|H^{1,0}(T))\) is a root of unity.
\end{lemma}

\begin{proof}
Put \(V=H^1(T,\Q)\) and \(D_\Q=\NS(T)_\Q\).  For each eigenvalue
\(\theta\) of \(A_\C\), let \(V_\theta\) be its generalized eigenspace and set
\(d_\theta=\dim V_\theta\), \(p_\theta=\dim(V_\theta\cap H^{1,0}(T))\).
Complex conjugation gives
\(p_\theta+p_{\overline\theta}=d_\theta=d_{\overline\theta}\).

For \(|\theta|\ne1\), project \(c\) to the block
\(V_\theta\wedge V_{\overline\theta}\), with \(\bigwedge^2V_\theta\) understood for
real \(\theta\).  This projection preserves \(D_{\overline\Q}\).  Indeed
the Chinese-remainder projectors onto the generalized spaces are
polynomials in \(A\) with algebraic coefficients.  Apply them to the
two tensor slots separately and then alternate.  The resulting maps are
algebraic linear combinations of rational Hodge maps on \(\bigwedge^2V\);
these preserve rational \((1,1)\) classes, hence \(D_\Q\) by Lefschetz
\((1,1)\).  Separate slot projections avoid any collision of products of
eigenvalues in degree two.

On this block \(A_2=\bigwedge^2 A\) has the single generalized
eigenvalue \(\theta\overline\theta=|\theta|^2\ne1\).  The inverse of \(A_2-1\)
there is a finite polynomial in its nilpotent part.  Project the congruence
for \(c\) to the block and apply this inverse.  It follows that its block
\(c_{\theta,\overline\theta}\) belongs to \(D_\C\).  Positivity says that this
block is a perfect tensor between \(V_\theta\) and \(V_{\overline\theta}\),
pairing their opposite Hodge types: in a basis of \(H^{1,0}\) the relevant
matrices are positive-definite principal Hermitian blocks.  For real
\(\theta\) the corresponding alternating tensor is nondegenerate.

The intersection of this block with \(D_\C\) is defined over
\(\overline\Q\).  Nondegeneracy is the nonvanishing of a determinant,
so it contains a nondegenerate tensor with algebraic coefficients.  For
any Galois automorphism \(\gamma\), the conjugate tensor remains in
\(D_{\overline\Q}\), hence is still of type \((1,1)\) for the given
Hodge decomposition, and is perfect between \(V_{\gamma(\theta)}\) and
\(V_{\gamma(\overline\theta)}\).  Opposite-type perfection yields
\begin{equation}\label{floor:torus-hodge-count}
 p_{\gamma(\theta)}+p_{\gamma(\overline\theta)}
       =d_\theta\qquad(|\theta|\ne1).
\end{equation}
This does not replace \(\gamma(\overline\theta)\) by
\(\overline{\gamma(\theta)}\), and assumes no Galois invariance of the Hodge
decomposition.

The top-form scalar is the algebraic integer
\(z=\prod_\theta \theta^{p_\theta}=\det(A|H^{1,0})\), also when generalized eigenspaces
are nontrivial.  For any Galois automorphism \(\delta\), reindex by the
image roots and pair complex conjugates:
\[
 \begin{aligned}
 \log|\delta z|
 &=\frac12\sum_\theta
       \bigl(p_{\delta^{-1}(\theta)}
             +p_{\delta^{-1}(\overline\theta)}\bigr)\log|\theta|\\
 &=\frac12\sum_\theta d_\theta\log|\theta|
   =\tfrac12\log|\det A|=0.
 \end{aligned}
\]
For \(|\theta|\ne1\) the second equality uses
Equation~\eqref{floor:torus-hodge-count} with \(\gamma=\delta^{-1}\); the other
terms vanish.  The last equality is unimodularity of the integral action.
Kronecker's theorem now proves the assertion.
\end{proof}

\begin{proof}[Proof of Proposition~\ref{floor:top-form-scalar}]
Apply Lemma~\ref{floor:finite-base}.  After a power, \(g\) acts trivially
on the canonical Iitaka base.  Use the modification \(\pi:U'\to U\) and
the map \(f:U'\to Y\) in its construction, and restrict over a general
point where the graph and the fibration are smooth and the fiber meets the
isomorphism locus of \(\pi\).  Dividing \(\alpha\) by a local base volume
gives a nonzero top form on the smooth compact fiber \(F\), which has
\(\kappa(F)=0\); its returning map scales the form by the corresponding
power of \(\mu\).  A point fiber gives scalar one.  For a positive-dimensional
fiber, let \(\omega_U\) be a K\"ahler form representing the class in the
proposition and put \(\beta_F=(\pi|_F)^*\omega_U\).  This is a smooth
semipositive form, strictly positive on the dense open where \(\pi\) is
an isomorphism.  Pulling the original graph congruence to \(U'\) and
restricting to the fiber shows that the two pullbacks of \([\beta_F]\) on
a common resolved graph of the returning map of \(F\) differ by real line classes.
This weaker positivity suffices below.

Since \(\dim F\leq d<n\), Corollary~\ref{prog:terminal-kappa-zero}, using
\(\Ggood_{\dim F}\), supplies a terminal compact K\"ahler minimal model
\(M\) of \(F\) with torsion canonical line.  The descended
nonzero reflexive top form trivializes \(K_M\) itself.  Indeed a power of
that form, in a trivialization of a multiple of \(K_M\), is a nonzero
holomorphic function on the compact connected \(M\), hence a nonzero
constant.  On a graph resolution of the returning map the two divisors of
this volume form agree; terminality makes their positive components exactly
the exceptional primes.  The returning map on \(M\) is therefore small
in both directions.

The K\"ahler Beauville--Bogomolov decomposition
\cite[Theorem~A]{BGL} supplies a finite cover, \'etale in codimension one,
of \(M\) of the form
\[
 P=T\times\prod_i H_i\times C,
\]
where \(T\) is a torus, the \(H_i\) are primitive symplectic factors,
and \(C\) is the product of the strict Calabi--Yau factors of dimension
at least three.  One-dimensional factors are elliptic curves and are
absorbed into \(T\); two-dimensional factors are placed among the
symplectic ones.  An absent factor is a point and contributes scalar one.
Reflexive forms on these klt spaces can be computed on resolutions by
\cite[Corollary~1.8]{KebekusSchnell2021}.  A power of the
small returning map lifts bimeromorphically to \(P\).  To justify this,
purity makes the restriction of the cover over \(M_{\rm reg}\) finite
\'etale \cite[Section~3.2, before Lemma~3.7]{BGL}.  This restriction is
connected and, after choosing base points, determines a finite-index subgroup
of \(\pi_1(M_{\rm reg})\), well-defined up to conjugacy.
Remove the codimension-two exceptional sets from the smooth loci.  Removing
an analytic subset of complex codimension at least two from a manifold does
not change its fundamental group.  The compact normal connected \(M\) is
reduced and paracompact, with bounded local tangent dimension.  Its analytic
pair \((M,M_{\rm sing})\) therefore has a real analytic embedding
\cite[Theorem~1]{AcquistapaceBrogliaTognoli1979}
and a compatible locally finite triangulation
\cite[Section~3, Theorem~1]{Lojasiewicz1964}.  Compactness makes the
triangulation finite; barycentric subdivision gives the smooth locus finite
CW homotopy type, so its fundamental group is finitely generated.
There are only finitely many subgroups of
the fixed index of the cover.  A power preserves the conjugacy class of its covering subgroup,
so it lifts on this open, and normalization over the proper graph extends
the lift bimeromorphically.  We may take further powers throughout.

We explain why the factors can be considered separately.  All holomorphic
one-forms on \(P\) come from \(T\).  Their pullbacks show on the smooth
open that the torus coordinate of the map depends only on the torus
coordinate.  Modulo the ideal generated by positive-degree torus forms,
the two-forms have basis the symplectic forms \(\eta_i\).  Write
\(\dim H_i=2e_i\).  The highest nonzero power of \(\sum a_i\eta_i\)
has order \(\sum_{a_i\ne0}e_i\); thus the locus where its full power
vanishes is the union of the coordinate hyperplanes.  The map permutes
these hyperplanes and hence the lines \(\C\eta_i\), matching their
nilpotence orders \(e_i+1\).  No torus two-form can be added to an image
of \(\eta_i\), since wedging that term with the \(e_i\)-th power of the
corresponding target symplectic form would
contradict that nilpotence order.  The kernels of the forms now show that
each symplectic coordinate depends only on the corresponding symplectic
factor.  Take a power to remove the permutations.

If \(C\) is positive-dimensional, it has no holomorphic form of degree
\(\dim C-1\): the
form algebra of a strict factor has only degrees zero and its top degree,
and no factor has dimension one.  In the pullback of the volume form of
\(C\), the component with one differential outside \(C\) and all the
others in \(C\) therefore vanishes.  The derivative in the \(C\)
directions has full rank, because the whole map is bimeromorphic and the
other coordinates have just been separated.  Its cofactors then force the
derivative of the \(C\) coordinate in outside directions to vanish.
Thus that coordinate also depends only on \(C\).  These conclusions on a
dense smooth open give factor bimeromorphic maps by taking their proper
graphs.  Their degrees are one because their product has degree one.

The scalar for \(C\) is a root of unity by projective pluricanonical
finiteness.  In detail, \(H^{2,0}\) of a smooth compact K\"ahler
resolution of \(C\) vanishes by the strict-factor form algebra and
reflexive extension.  Rational approximation of a K\"ahler class and
Kodaira's theorem make this resolution projective.  Thus \(C\) is
Moishezon; it is projective by the K\"ahler projectivity criterion for
rational singularities \cite[Corollary~6']{Namikawa2002}.  Its canonical
line is trivial and its returning comparison is crepant, as seen from the
volume form.  Chow's theorem algebraizes its proper analytic graph.  The
projective log pluricanonical representation theorem
\cite[Theorem~1.1]{FujinoGongyo2014} gives finite scalar image.

We spell out how the polarizing class reaches the other factors without
asserting a cohomology splitting for the singular product \(P\).  Choose
projective resolutions of its factors with smooth compact K\"ahler sources,
and form the smooth compact
K\"ahler product
\[
 X_P=T\times\prod_i\widetilde H_i\times\widetilde C.
\]
Resolve the map \(P\dashrightarrow F\) together with \(X_P\to P\).  This
gives a smooth compact K\"ahler space \(R\), a modification \(v:R\to X_P\),
and a morphism \(r:R\to F\).  Put \(\beta=r^*\beta_F\).  It is smooth
semipositive and strictly positive on a dense open, since \(r\) is
generically finite.  Put \(a=v_*[\beta]\).  The smooth
modification formula gives \([\beta]=v^*a+\{E_v\}\) for an exceptional real
divisor \(E_v\).  The graph congruence for \([\beta_F]\) pulls to the cover and,
by Gysin pushforward, gives the graph congruence for \(a\) modulo real
line classes; the exceptional differences are themselves line classes.
Pass to a common refinement and use the product of resolutions of the
separated factor graphs to compute this congruence; refinements change it
only by exceptional line classes.  The restriction \(a_i\) of \(a\) to one factor slice
is independent of the other coordinates, by K\"unneth on the smooth
product \(X_P\).  Restricting the product graph therefore compares
\(a_i\) to itself modulo line classes on the factor graph.

Choose the factor slice generally so that it meets the open where \(\beta\)
is strictly positive and meets each \(v\)-exceptional center in codimension at least two,
or avoids that center.  This is possible by the fiber dimension theorem
applied to each center and the projection to the other factors.  On the
resolved strict transform of this slice, the restriction \(\beta_i\)
is still semipositive and strictly positive on an open, and the restriction
of \(E_v\) is exceptional over the factor.  Thus the pullback of \(a_i\)
differs from \([\beta_i]\) only by the class of an exceptional real divisor. This
gives both the required positivity and the congruence on each smooth
factor model.

For a symplectic factor, its symplectic volume trivializes its canonical
line.  Canonical singularities preserve the canonical ring on resolution,
so a smooth resolution has Kodaira dimension zero.  Its dimension is
\(\dim H_i\leq\dim F<n\).  Corollary~\ref{prog:terminal-kappa-zero}, using
\(\Ggood_{\dim H_i}\), runs from that resolution to a terminal model.  It is still primitive
symplectic: the unique two-form extends, and its top power trivializes the
canonical line and is nondegenerate on the smooth locus; irregularity and
the form algebra are unchanged.  Take a common projective resolution of
this terminal model carrying the pullback of \(\beta_i\).
Lemma~\ref{prog:cohomology-transport} decomposes its class as the pullback
of a class \(c\) on the terminal model plus the class of an exceptional real divisor.
The line-trace assertion of that lemma sends the line Chern classes in
the factor graph congruence to rational line classes on the terminal model.
Pushing through the exceptional quotient therefore shows that the returning
graph compares \(c\) to itself modulo line classes.  These are the
hypotheses of Lemma~\ref{floor:symplectic-scalar}.
That lemma makes the scalar on its two-form, and therefore on its top form,
a root of unity.

For the torus, push the semipositive form to the torus as a positive current
and average it over translations.  The average is a smooth invariant
positive Hermitian representative: for every nonzero tangent direction its
average is strictly positive because the original form is strictly positive
on an open.  A bimeromorphic map of a torus is an affine holomorphic
automorphism; its linear part preserves the integral \(H^1\) lattice.
The comparison modulo \(\NS_\R\) and Lemma~\ref{floor:torus-scalar}
make its top-form scalar a root of unity.  The product of the factor
scalars is the appropriate power of \(\mu\).  It is a root of unity,
and hence so is \(\mu\).
\end{proof}

\subsubsection{Finite image on vertical strata}

\begin{proof}[Proof of Proposition~\ref{floor:finite-action}]
Put \(d=\dim Z\) and \(m=qk\), so sections of \(L_Z^k\) are meromorphic
\(m\)-pluricanonical forms.  Let
\[
 E_m=\left\{s\in H^0(Z,mJ_Z):\int_Z|s|^{2/m}<\infty\right\}.
\]
The integral is computed on the regular locus, or on any resolution by
change of variables.  This is a vector space: for \(2/m\leq1\) the
elementary power inequality preserves integrability under sums.  On a log
resolution an adjoint section is locally
\[
 h(z)\frac{(dz_1\wedge\cdots\wedge dz_d)^{\otimes m}}
                 {\prod_i z_i^{mb_i}},\qquad b_i\leq1.
\]
It is integrable if \(h\) vanishes on every unit component; smaller
coefficients cause no obstruction.  Every unit valuation has center in
\(C_Z\).  Thus \(E_m\) contains pullbacks of base sections vanishing
along \(f_Z(C_Z)\).  For large \(k\) these define \(Y_Z\) birationally:
multiply one nonzero section of a high ample power vanishing on that proper
subset by a complete very ample system.  If \(Y_Z\) is a point,
verticality forces \(C_Z=\varnothing\) and a nonzero section is
integrable.  In particular \(E_m\ne0\).

The function \(s\mapsto\int|s|^{2/m}\) is continuous on \(E_m\) and
positive away from zero.  In an integrable basis, the densities of bounded
linear combinations are dominated by a constant times the sum of the basis
densities, proving continuity by dominated convergence.  Its invariance
under arrows makes the action and its inverse uniformly bounded in any
norm on \(E_m\).  Moreover this representation detects the representation
on all of \(H^0(Z,mJ_Z)\).  If an arrow acts identically on \(E_m\), it
acts identically on \(Y_Z\) because that system is birational.  For
\(0\ne s\in E_m\) and any other section \(t\), the ratio \(t/s\) comes
from the rational function field of \(Y_Z\).  The arrow fixes both this
ratio and \(s\), and hence fixes \(t\).

Take an eigenform \(s\in E_m\setminus\{0\}\) for the transport
\(\gamma_*\) of one self-arrow, with eigenvalue \(\lambda\).  In the
orientation fixed above, \(\gamma_*=(\gamma^{-1})^*\); put
\(g=\gamma^{-1}\), so \(g^*s=\lambda s\).  Form the full normalized analytic cyclic root
cover of this meromorphic pluriform and take a projective resolution of
all its components.  Locally, if \(s=a\eta^{\otimes m}\) in a
meromorphic canonical frame, the cover is the normalization of
\(t^m=a\), and its tautological top form is \(\alpha=t\eta\).
These descriptions glue when the frame changes.  The finite analytic
normalization and this full-root construction are also described in the
proof of \cite[Lemma~7.1]{BL}.  On the smooth compact K\"ahler resolution
\(U_s\) of the full cover, let \(\pi_s:U_s\to Z\) be the composite map.
Change of variables gives
\[
 \int_{U_s}|\alpha|^2=(\deg\pi_s)\int_Z|s|^{2/m}<\infty.
\]
A meromorphic top form with finite local \(L^2\) norm has no pole, by
the one-variable integral transverse to a putative pole.  Thus \(\alpha\)
is holomorphic.  Choosing \(\mu^m=\lambda\) lifts \(g\)
bimeromorphically to the full cover, possibly permuting components, with
\(\widetilde g^*\alpha=\mu\alpha\).

The lift preserves a K\"ahler class modulo \(\NS_\R\).  A sufficiently
large multiple of the pullback of \(c_Z\), plus a relatively ample line
class for the finite cover and its projective resolution, is K\"ahler on
\(U_s\).  Pulling up Equation~\eqref{floor:class-comparison} compares the first
summand, and the changes of the relatively ample classes are line classes.
After a power fixes one connected component, that component has dimension
\(d=\dim Z<n\), so Proposition~\ref{floor:top-form-scalar} applies there.
It follows that
\(\mu\), and hence \(\lambda\), is a root of unity.

The possible orders are bounded in this fixed degree \(m\).  We give the
parameter argument because pointwise torsion alone would not make a bounded
group finite.  Fix a smooth resolution \(X\) of \(Z\) and an effective
divisor \(D\) clearing the poles of a basis of \(E_m\).  Put
\(Q_m=\PP(E_m^\vee)\), the projective parameter space of lines of forms
under our quotient convention.  The universal section on \(X\times Q_m\)
belongs to
\[
 (K_X(D))^{\otimes m}\boxtimes\mathcal O_{Q_m}(1).
\]
Pull the parameter space back by the finite coordinate-power map
\([z_0:\cdots:z_b]\mapsto[z_0^m:\cdots:z_b^m]\).  Its pullback of
\(\mathcal O(1)\) is \(\mathcal O(m)\), so the displayed line now has
an \(m\)-th root \(\mathscr L\).  The cyclic algebra
\(\bigoplus_{i=0}^{m-1}\mathscr L^{-i}\), with multiplication defined
by the universal section, defines a finite locally free family
\(\mathcal C\) of rank \(m\).
Every parameter is a nonzero form.  Each fiber is therefore reduced:
its algebra is torsion-free over the smooth \(X\) and is generically
squarefree.  Every fiber component meets the inverse image
\(W\subset\mathcal C\) of the nonvanishing locus of the universal
section; there the cover is finite \'etale.  Normalize \(\mathcal C\)
and resolve the normalization projectively.  The composite
\(\rho:\widetilde{\mathcal C}\to\mathcal C\) can be chosen to be an isomorphism
over \(W\).  Put \(E_\rho=\rho^{-1}(\mathcal C\setminus W)\).
The locus
\(\{t:\dim(E_\rho)_t\ge d\}\) is proper analytic by the proper
fiber-dimension theorem: every component of the total family meets
\(W\), and its general fiber has dimension \(d\).  Remove this locus
and the critical values of the resolved map.  Over the remaining dense
open the family is smooth and proper, and each fiber is the disjoint
union of smooth resolutions of the corresponding root-cover components,
because every component meets the locus where \(\rho\) is an isomorphism.
Proper smooth transport makes the sum of their middle Betti numbers
constant there.  Resolve
the finitely many irreducible components of the proper analytic complement.
On each such resolution pull back the original finite cyclic family, take
its reduction, and normalize and resolve anew; restricting the previous
normalization is not required.  The parameter dimension decreases,
so finitely many such steps cover all parameters.  There is consequently a
uniform bound \(B_m\) for the middle Betti number of the entire disjoint
union of selected smooth resolutions of every full cover.  This does not assert a bound for
arbitrary further resolutions.

On the disjoint union of the selected resolutions, let \(p_1,p_2\) be the
projections of a smooth resolved graph of the lift.  The integral Gysin
endomorphism \((p_1)_!p_2^*\) on middle cohomology has \(\mu\) as an eigenvalue:
the nonzero holomorphic top form has a nonzero cohomology class,
\(p_2^*[\alpha]=\mu p_1^*[\alpha]\), and \((p_1)_!p_1^*=\id\).
If the order of \(\mu\) is \(a\), its
cyclotomic minimal polynomial therefore has degree \(\varphi(a)\leq B_m\).
There are only finitely many such \(a\).  Thus the eigenvalues
\(\lambda=\mu^m\) on \(E_m\) lie in a fixed finite set.

The bounded group on \(E_m\) has compact closure in \(\mathrm{GL}(E_m)\).
Every element of that closure still has its eigenvalues in this finite set,
by continuity of characteristic polynomials.  On the identity component
the characteristic polynomial must be that of the identity.  A compact
linear group is unitarizable, so an element all of whose eigenvalues are
one is the identity.  The identity component is trivial; a compact Lie
group with this property is finite.  The image on \(E_m\), and therefore
on the full section space, is finite.
\end{proof}

\subsection{Compatible sections on the whole boundary}

We now use restriction, transport, and finite self-arrow image to construct
sections simultaneously on all boundary strata.

For \(d\geq0\) and \(k\geq1\), a system of tuples in degree \(k\) through dimension
\(d\) means a vector subspace
\[
 \mathcal V_d(k)\subseteq
       \bigoplus_{\substack{Z\subset S\text{ a stratum}\\\dim Z\leq d}}
                      H^0(Z,L_Z^k).
\]
It is \emph{compatible} if every tuple agrees under the residue restrictions
for every incidence, and \emph{invariant} if every arrow carries its source
component to its target component.  It \emph{generates} if for every point
of every stratum some tuple has nonzero value there.  Compatibility and
invariance are linear conditions.  If such a system generates, the span
of componentwise \(a\)-th powers of its tuples gives a compatible
invariant generating system in degree \(ak\).  Thus we may always pass to
a sufficiently large divisible later degree.

\begin{proposition}[Generating tuples]\label{floor:tuples}
There is a degree \(k>0\) and a compatible invariant generating system
of tuples through dimension \(n-1\).
\end{proposition}

\begin{proof}
At the point strata use the same scalar in the canonical zero-form
generator \(1\).  The even residue convention identifies these generators
along every path and under every point arrow.  This gives the required
system in dimension zero; if there are no points the empty system is
understood.

Suppose the system has been constructed through dimension \(d-1\).
Replace it by powers and their span in a common degree \(k\) chosen so that
Lemma~\ref{floor:restriction} applies on
every \(d\)-stratum, Proposition~\ref{floor:finite-action} applies on every
vertical \(d\)-stratum, and
\(\mathcal I_{f_Z(C_Z)}\otimes N_Z^k\) is generated for every
\(d\)-stratum \(Z\).  These conditions hold in a common divisible tail:
there are finitely many strata, and the last condition is Serre's theorem.
For a \(d\)-stratum \(Z\), Lemma~\ref{floor:adjunction} glues
the lower tuple to a section on its whole \(C_Z\).  Its residues satisfy
the one link in Lemma~\ref{floor:restriction}, if that link is present,
because the lower tuple is invariant.  That lemma extends the section to
\(Z\).

Let \(\mathcal P_d(k)\) be the vector space of all tuples through
dimension \(d\) whose lower part is in \(\mathcal V_{d-1}(k)\) and whose
new components restrict to that lower part.  Call its elements
\emph{pre-tuples}.  Its projection to
\(\mathcal V_{d-1}(k)\) is surjective, since the finitely many extensions
can be chosen independently.  This system already generates.  To check a
point \(z\in Z\), put \(y=f_Z(z)\).  If \(y\in f_Z(C_Z)\), choose a
floor point over \(y\) and a lower tuple nonzero there.  Every extension is
the pullback of a section of \(N_Z^k\), so it is nonzero at every point
over \(y\), including \(z\).  If \(y\notin f_Z(C_Z)\), the chosen generation
of \(\mathcal I_{f_Z(C_Z)}\otimes N_Z^k\) supplies a base section nonzero
at \(y\).  It pulls back to a section
vanishing on \(C_Z\); combine it with the zero lower tuple and zero
components on the other new strata.  This also treats an empty floor.
Surjectivity to the lower system preserves generation at lower points.

We impose invariance in dimension \(d\).  For a dominating floor,
injectivity of restriction and Lemma~\ref{floor:transport} already make
every pre-tuple invariant on that stratum.
Verticality is invariant under arrows, so the remaining orbits consist
entirely of vertical strata.

For each such orbit choose a representative \(Z\), and let \(G_Z\) be
the finite isotropy image on \(H^0(Z,L_Z^k)\) given by
Proposition~\ref{floor:finite-action}.  For a pre-tuple with component
\(s_Z\), form the norm
\[
 P_Z(s_Z)=\prod_{g\in G_Z}g(s_Z)
       \in H^0(Z,L_Z^{k|G_Z|}).
\]
Choose a common positive integer \(a\) divisible by all \(|G_Z|\), use
\(P_Z(s_Z)^{a/|G_Z|}\), and transport it to each member of that orbit.
This is independent of the chosen arrow: changing that arrow by an
isotropy arrow merely permutes the factors in degree \(k\).  No finiteness
assertion in degree \(ak\) is needed.  On strata with dominating floor
use \(s_Z^a\), and on all lower strata use the componentwise \(a\)-th
power of the assigned tuple.

The resulting tuple is compatible.  By Lemma~\ref{floor:transport},
every transported factor in a norm has exactly the prescribed lower
restriction.  The norm and its indicated power therefore restrict to
the \(a\)-th power of that restriction, the same value used on all lower
strata and in all other orbits.  It is invariant in dimension \(d\) by
construction, and remains invariant below it.

These tuples still generate.  Fix a point \(z\) in an orbit member.
For each of the finitely many transported factors, choose a point over
\(z\) on a resolution of its comparison graph.  Equality of the
pulled-back invertible adjoint lines identifies the fiber values.  The
factor is nonzero at \(z\) exactly when the representative component
\(s_Z\) is nonzero at the corresponding point of \(Z\).  Evaluation at each
such point is a nonzero linear functional on the generating space
\(\mathcal P_d(k)\).
A finite union of their proper kernels cannot cover a complex vector
space.  One pre-tuple makes all factors nonzero and gives a norm nonzero
at \(z\).  At a lower point a nonzero assigned value stays nonzero after
its \(a\)-th power.  Taking the linear span of all the constructed tuples
therefore gives a compatible invariant generating system in degree \(ak\).
This completes the induction on \(d\).
\end{proof}

\begin{proof}[Proof of Theorem~\ref{floor:generation}]
Take the system in Proposition~\ref{floor:tuples}.  Its components on the
prime components of \(S\) agree on every common lower stratum, so
Lemma~\ref{floor:adjunction} descends each tuple uniquely to a section of
the actual line \(\mathcal O_V(qkJ)|_S\).  At any \(x\in S\), choose a
prime component through \(x\) and a tuple nonzero there.  It is the pullback
of the descended value in the same invertible line fiber at \(x\), so
that value is nonzero.  The finite-dimensional span of these descended
sections generates at every point of the reduced \(S\).  This is the
claimed global generation of an actual Cartier multiple.
\end{proof}
